\begin{lemma}[Situation 7.3, Claim \(j-1\): empty case]
Fix a rung-local minimum hierarchy counterexample \(C\), set

\[
 j=\noderef{HierarchyPathNumber}(C.\mathit{hierarchy}),\qquad
 \ell=\noderef{HierarchyFullTailLength}(C),
\]

and assume \(2\le j<\ell\).  Assume every vertex has a missing color
under every proper exactly-one-uncolored replacement coloring.  If
\(\noderef{HierarchyEligibleCommonMissingIndices}\) is empty, then there
nevertheless are an actual
\noderef{HierarchyReplacementPresentation} \(P\) and a color \(\alpha\)
missing under \(P.\mathit{coloring}\) both at the terminal vertex and
somewhere in the exact prefix \(T(y_{j-2})\), such that, with the notation
\(D=D_{n,q}\), \(\Gamma^q=\bigcup_{\eta\in D}\Gamma^q_\eta\), and
\(M_0=\overline\varphi(P.\mathit{replacementBase}.\mathit{core}-T_n)\),
either

\[
                       \alpha\notin\Gamma^q\cup M_0,          \tag{A}
\]

or \(\alpha\in\Gamma^q_\eta\) for some \(\eta\in D\) whose first
missing vertex occurs no later than \(y_{j-2}\).
\end{lemma}
\begin{proof}
We first construct the identity presentation \(P_0\).  Take the original
coloring, the reflexive
\noderef{CorrespondingExtendedTashkinovTree}, the maximum-property field
of \(C\), the original \noderef{HierarchyBase}, \(C\) itself, the
reflexive \noderef{SameGoodHierarchy}, and the supplied
\noderef{HierarchyMinimumCounterexampleContext}.  These are exactly all
fields of \noderef{HierarchyReplacementPresentation}; thus \(P_0\)
belongs to the replacement family, without an existence assumption.
Write \(H=C.\mathit{hierarchy}\), let \(q\) be its last stored index, and
put

\[
\begin{gathered}
 A=\noderef{HierarchyStarredStage}(H,q),\qquad
 B=T(y_{j-2}),\qquad B^+=T(y_{j-1}),\\
 B_{\rm pen}=T(y_{\ell-1}),\qquad
 D=\noderef{HierarchyStoredColorsAt}(H,q),\\
 M_0=
 ((V(\mathit{base.core})\setminus V(T_n))
       \mathbin{\mathop{\bigcup}}\overline\varphi(v)),\qquad
 \Gamma^q=\bigcup_{\eta\in D}C.\mathit{good}.\Gamma^q_\eta .
                                                               \tag{1}
\end{gathered}
\]

All prefixes in (1) are the exact
\noderef{HierarchyFullTailPrefix}es.  Apply
\noderef{HierarchyCounterexamplePenultimateElementary} to the properness,
exactly-one-uncoloredness, complete lower-rung package, fewer-divider
induction, lexicographic minimality, and \(2n+11\) fields of the supplied
minimum context.  It gives

\[
              V(B_{\rm pen})\text{ is elementary under }\varphi. \tag{2}
\]

Apply \noderef{HierarchyCounterexampleTerminalCommonMissing} to those
same six fields.  Choose

\[
 \alpha\in\noderef{TreeMissingColors}(\varphi,B_{\rm pen})
       \cap\noderef{MissingColors}(G,\varphi,y_\ell).          \tag{3}
\]

The assumed emptiness of the eligible-index set says that no displayed
tail vertex \(y_t\), \(1\le t\le\ell-1\), can be an
\(\alpha\)-missing witness in (3), since \(P_0,t,\alpha\) would be a
member of the defining filter.  Every vertex of \(B_{\rm pen}\) is
either in \(A\) or is one of those tail vertices, by the entrywise prefix
definition and \noderef{TreeSequence}.  Hence

\[
                   \alpha\in\noderef{TreeMissingColors}(\varphi,A).
                                                               \tag{4}
\]

In particular \(\alpha\in\overline\varphi(B)\), because
\(A\preccurlyeq B\) and \(j\ge2\).  If
\(\alpha\notin\Gamma^q\cup M_0\), then \(P_0,\alpha\) already satisfy the
conclusion.  We therefore split the remaining possibilities exactly as
in the paper.

\textbf{Base-extra case.}
Assume

\[
                      \alpha\in M_0\setminus\Gamma^q.          \tag{5}
\]

The terminal-ladder missing set has cardinality at least \(2n+11\) by
the large-missing-set field of the minimum context.  Each stored color
accounts for one index and, by the two-colors and disjoint-pairs fields
of \noderef{GoodHierarchy}, \(\Gamma^q\) accounts for at most
\(2|D|\le2n\) colors.  Elementarity (2) identifies the missing colors
contributed by the base core and prevents double counting.  Deleting
\(M_0\) and \(\Gamma^q\) therefore leaves at least eleven colors missing
on \(A\); choose

\[
 \beta\in\noderef{TreeMissingColors}(\varphi,A)
                    \setminus(M_0\cup\Gamma^q).               \tag{6}
\]

The first \(\alpha\)- and \(\beta\)-missing witnesses lie in the
elementary prefix.  The appropriate case of \noderef{ChangeLemma}
identifies their component, and the terminal component through
\(y_\ell\) is distinct: equality would give three distinct missing
endpoints in one component, contrary to
\noderef{MissingColorIncidenceEndpoint},
\noderef{TwoColorComponentSpanningPath}, and
\noderef{TwoColorIncidenceDegreeAtMostTwo}.  Use
\noderef{TwoColorComponentsDisjoint}, switch the terminal component by
\noderef{TwoColorComponentSwap}, and call the coloring \(\sigma\).  It
is proper by \noderef{TwoColorComponentSwapProper} and leaves exactly
\(e_0\) uncolored because the switch fixes the uncolored value.

We verify the replacement rather than merely naming it.  For a vertex
\(v\in V(B_{\rm pen})\) outside the switched component, adjacency in
\noderef{TwoColorIncidenceGraph} fixes every affected incident edge.
For \(v\) inside it, disjointness from the first-witness component and
elementarity (2) make both affected colors present at \(v\).  Properness
gives one incident edge of each color; the switch exchanges those two
and fixes every other incident color.  Thus every missing set on
\(B_{\rm pen}\) is unchanged pointwise.  This is valid even when the
switched component meets the prefix internally.

Inspect the actual final \noderef{TashkinovSeriesStep} of the original
ETT and use \noderef{ExtendedTashkinovTree}.\mathit{belowFinal}.  In PE,
retain its terminal-ladder closure route and the exact parallel core and
closure from \noderef{HierarchyBase}.  In SE/RE, retain the individual
connecting edge, singleton connecting set, literal
\noderef{TreeSequenceAdjoinEdge} seed, and its route.  At positive
divider count \noderef{HierarchyNonparallelFirstEdge} places the seed
below \(A\); at divider zero the seed remains the route source.  Restrict
at \(B_{\rm pen}\) using
\noderef{TashkinovAugmentationSequenceRestrict}, transport each exact
piece using the pointwise missing-set equality, and concatenate in
\noderef{TashkinovAugmentationSequence}.  This verifies all route,
first-position, and Gamma-cutoff premises of
\noderef{StableChangeCasesOneTwoSplitTail}, separately for PE,
divider-zero SE/RE, and positive-divider SE/RE.  It returns a
corresponding same-rung maximum ETT on \(T\), rebuilt base, hierarchy,
goodness, and same-good data.

The switch makes \(\beta\) missing at \(y_\ell\) while its witness in
\(A\subseteq B\) remains missing.  Package the returned data as \(P_1\).
Its complete minimum context is transported field by field:
correspondence transports the lower-rung package and large estimate;
same-good equality transports divider count and the fewer-divider field;
the identical path number, stages, and terminal excess identify both
coordinates of \noderef{HierarchyCounterexampleMeasure}; and properness,
one-uncoloredness, and lexicographic minimality have just been verified.
By (6), \(P_1,\beta\) satisfy (A) and prove the claim in the base-extra
case.

\textbf{Gamma-pair case.}
It remains that
\[
 \alpha\in C.\mathit{good}.\Gamma^q_{\eta_m}
                    \quad\text{for some }\eta_m\in D.         \tag{7}
\]
We distinguish whether \(\eta_m\) is missing on \(B_{\rm pen}\).

First suppose it is not.  The avoidance field of
\noderef{GoodHierarchy} excludes \(\alpha\) from every displayed tail
edge through \(B_{\rm pen}\).  Apply \noderef{NineN} at the exact prefix
\(T(y_{\ell-2})\).  The complement alternative gives a color
\(\beta\) missing on that prefix, outside
\(M_0\cup\varphi\langle T-A\rangle\cup(\Gamma^q\cup D)\); the seven-color
alternative gives
\[
 \beta=\eta_h\in D\cap
     \noderef{TreeMissingColors}(\varphi,T(y_{\ell-2}))
\]
with
\((C.\mathit{good}.\Gamma^q_{\eta_h}\cup\{\eta_h\})\) disjoint from the
entire two-edge suffix \(T-T(y_{\ell-2})\).  The latter choice is possible
because the seven indexed three-color sets are pairwise separated by
the injectivity and disjoint-pairs clauses, while two exact edges use
only two colors.  In either alternative \(\beta\) is missing on
\(T(y_{\ell-2})\) and neither \(\alpha\) nor \(\beta\) is used on the
relevant tail interval.

Apply \noderef{ChangeLemma} to \((\alpha,\beta)\), separate its
first-witness component from the terminal component by the
three-endpoint argument, and switch the terminal component.  Use the
pointwise outside-versus-internal-vertex proof and the actual terminal
certificate construction just given, now at the appropriate split
prefix, to invoke \noderef{StableChangeCasesOneTwoSplitTail}.  The
result is an actual replacement presentation \(P_2\) with its own full
minimum context; \(\beta\) is missing both at \(y_\ell\) and at its
unchanged first witness \(v_\beta\).

Compare that exact first position with \(y_{j-2},y_{j-1}\), and the path
suffix.  If \(v_\beta\preceq y_{j-2}\), then the color-selection
alternative gives exactly (A), or the stated Gamma-pair alternative,
for \(P_2,\beta\).  If \(v_\beta=y_{j-1}\), then \(j-1\) is a member of
\noderef{HierarchyEligibleCommonMissingIndices}; apply
\noderef{HierarchySituationThreeClaimJMinusOneNonempty} to that now
nonempty finite family.  If \(y_j\preceq v_\beta\), its numerical index
is an eligible member at least \(j\), contradicting
\noderef{HierarchySituationThreeEligibleIndexBound}.  These alternatives
are exhaustive by the injective total order of the displayed vertices.

Finally suppose
\(\eta_m\in\noderef{TreeMissingColors}(\varphi,B_{\rm pen})\).
Because \(\eta_m\in D\), it is not missing on \(A\); hence its first
tail occurrence is a vertex \(y_t\) with \(1\le t\le\ell-1\).  If
\(t\le j-2\), then the original \(P_0,\alpha\) satisfy the Gamma-pair
alternative in the conclusion: (4) puts \(\alpha\) in \(B\), (3) puts
it at the terminal vertex, and the first \(\eta_m\)-position is at most
\(|B|\).

Assume \(t\ge j-1\).  Both \(\alpha\) and \(\eta_m\) are unused on the
displayed segment through \(y_t\), by Gamma avoidance and first
missingness.  There is one and only one possible exception: if \(q=0\),
the last extension is SE or RE, \(A=T_n\), and the first exact displayed
edge is the connecting seed colored \(\eta_m\in D_n\).  In that case use
the individual seed and the terminal-stability conclusion of
\noderef{StableChangeCasesOneTwo}; it proves that this edge lies outside
the terminal component and retains its value.  Thus the exception is
handled without claiming generic tail-edge invariance.

Apply \noderef{ChangeLemma} to \((\alpha,\eta_m)\), separate the
first-witness component from the terminal component, switch the latter,
and use the pointwise missing-set proof plus the applicable actual
terminal certificate to construct a full replacement presentation
\(P_3\).  The same fieldwise correspondence and same-good compositions
give \(P_3.\mathit{minimumContext}\), and \(\eta_m\) is missing under
\(P_3\) at both \(y_t\) and \(y_\ell\).  If \(t=j-1\), invoke
\noderef{HierarchySituationThreeClaimJMinusOneNonempty}; if \(t\ge j\),
invoke \noderef{HierarchySituationThreeEligibleIndexBound} for a
contradiction.  Together with the already completed case \(t\le j-2\),
this exhausts the Gamma-pair case and proves the empty-case claim.
\end{proof}
